\chapter{Workspace characterisation and its consequences}
\label{ch:workspace}

\Cref{ch:resultsmain} gives the headline finding and the mount fix it
produced. This appendix is the chain of measurements behind it.

\section{The two arms share no reachable workspace}

Over a 63-cell grid spanning the frontal working volume, arms verified at
home, no cell was reachable by both arms (\cref{fig:disjoint} in
\cref{ch:resultsmain}). The result survives freeing the wrist yaw (raising
each arm's individual reach from 16/18 to 21/21 cells, still 0/63 shared)
and additional solver retries. The right arm reaches
$x\le\SI{-0.10}{\metre}$, the left $x\ge\SI{0.15}{\metre}$: a
\SI{0.25}{\metre} dead band, measured at $y=\SI{0.35}{\metre}$,
$z=\SI{1.10}{\metre}$, yaw free, three repeats per point.

This holds within $x\in[-0.6,0.6]$, $y\in[0,0.4]$, $z\in[0.90,1.30]$
\si{\metre}, with the as-built mount and the legacy home joint angles -- the
arms share no workspace \emph{as mounted and as parked}, and the cause is
not the mount's rotation (proved invariant) but the arms' asymmetric home
joint angles, which are ground truth because the simulation-to-hardware
bridge replays them onto the real arm.

\subsection{The fix is buildable}

A parametric sweep moving the mounts inboard, forward and down, with
additional forward tilt, raises shared cells from 3 to \textbf{22 of 63}
within a buildable \SI{0.20}{\metre} mount separation (splaying the mounts
outward \emph{hurts}). Recorded and deliberately not applied at the time
(changing the mount invalidates the home pose, the anchor and every
clearance figure), then applied later as one deliberate pass with the whole
task set re-measured:

\begin{table}[htbp]
  \centering
  \caption[Before and after the mount change]{}
  \label{tab:mountapplied}
  \begin{tabular}{@{}lll@{}}
    \toprule
    Quantity & Before & After \\
    \midrule
    Clearance at home & \SI{0.1610}{\metre} & \SI{0.2202}{\metre} \\
    Nearest robot part to wearer at home & the mount (no joint moves it) & the moving chain \\
    Task objects resting on the surface & not achievable & achieved \\
    Pick-and-place targets off centre & \SI{0.530}{}/\SI{0.350}{\metre} & \SI{0.290}{\metre} either side \\
    \bottomrule
  \end{tabular}
\end{table}

\section{One boundary, two independent constraints}

Wearer clearance and top-down grasp feasibility, measured for different
reasons, land on the same lateral boundary: clearance rises monotonically
and crosses the \SI{150}{\milli\metre} floor at \SI{0.25}{\metre}; top-down
grasp candidates fail entirely (0/9) at $|x|=\SI{0.25}{\metre}$ and succeed
entirely (9/9) at $|x|=\SI{0.30}{\metre}$, with collision checking disabled
too -- kinematic infeasibility of the wrist orientation, not obstruction.

\section{What the pinned wrist costs at the work point}
\label{sec:work-point}

Reach under each orientation policy, walked from home, understates the cost
badly because home is in free space. Walked instead from
$(0.4,0.175,\SI{1.12}{\metre})$ on the working surface, with the
\SI{150}{\milli\metre} floor enforced throughout:

\begin{table}[htbp]
  \centering
  \small
  \caption[Reach from home and from the work point]{Mean reach under each
  orientation policy. The two columns disagree by a factor of seven, because
  reach is meaningless without the point it was measured from.}
  \label{tab:work-point}
  \begin{tabular}{@{}lcc@{}}
    \toprule
    Orientation policy & from home (\si{\metre}) & from work point (\si{\metre}) \\
    \midrule
    pinned to the anchor (delivered) & \num{0.460} & \textbf{\num{0.065}} \\
    within a \SI{15}{\degree} cone   & \num{0.492} & \num{0.304} \\
    within a \SI{45}{\degree} cone   & \num{0.530} & \num{0.325} \\
    position only                    & \num{0.547} & \num{0.363} \\
    \bottomrule
  \end{tabular}
\end{table}

At the point where the work happens, the delivered policy leaves the arm
\SI{65}{\milli\metre} of mean reach against \SI{363}{\milli\metre} with the
wrist free -- a factor of \num{5.6}. The mechanism is redundancy, not the
wearer: no walk from the work point stops on the clearance floor, they stop
because there is no solution -- a fixed six-degree-of-freedom pose on a
seven-degree-of-freedom arm spends the entire redundancy, and pinning the
wrist removes the joint that would otherwise move the elbow clear while the
hand stays put. Freeing the roll alone buys \SI{0}{\milli\metre} at the work
point; almost all recovery comes from the first \SI{15}{\degree}
(\SI{80}{\percent} of the pinned-to-free gap), which the follower's own
\SI{15}{\degree} solving cone already supplies -- the reason the task set
solves at all.

\section{The centre of the working space, asked four ways}
\label{sec:centre}

\begin{table}[htbp]
  \centering
  \small
  \caption[The centre, four measurements]{Four attempts to place usable work
  within \SI{100}{\milli\metre} of the body centre line, at ten repeats.}
  \label{tab:centre}
  \begin{tabular}{@{}p{0.36\linewidth}p{0.56\linewidth}@{}}
    \toprule
    Question & Answer \\
    \midrule
    Innermost workable column & $|x|=\SI{0.325}{\metre}$ (L), \SI{0.450}{\metre} (R) \\
    Same, wearer's arms deleted entirely & \SI{0.300}{}/\SI{0.375}{\metre}; every column
      \SIrange{0.125}{0.300}{\metre} binds on the \textbf{torso} \\
    Objects resting on a surface, \num{3360} configurations & \textbf{no
      configuration} within \SI{100}{\milli\metre} of centre \\
    Same, top-down approach & \textbf{0 of 840} cells, any height \SI{0.70}{}-\SI{1.10}{\metre} \\
    \bottomrule
  \end{tabular}
\end{table}

The constraint is the torso, not the wearer's limbs: sweeping five arm
postures moves the boundary by at most \SI{75}{\milli\metre}, and two
postures (arms clasped behind the back, arms out to the sides) place the
wearer's own limbs inside the clearance floor at home, before the robot
does anything. Neither arm crosses the centre line at all -- zero of ten
solutions, both arms, every cross-side slot -- which is why the
pick-and-place task draws the \emph{side} and lets colour follow it.

\section{Consequences for the task set}

Inter-arm handover is impossible: no control, planning or autonomy change
can produce a point neither arm reaches. A bimanual task must span the dead
band or exploit simultaneity. The coupled carrier's original
\SI{310}{\milli\metre} width placed each grip inside the dead band (10 of 25
waypoints failed once the arms were verified at home); widened to
\SI{500}{\milli\metre} it satisfies both boundaries, and the carrier's sag
tolerance was re-derived to keep the failure threshold inside the reachable
band rather than outside it, where it would measure nothing.
